\section{Introduction}
\label{sec:intro}

Critical infrastructures---power dispatch, datacenter scheduling, financial settlement---are increasingly automated by tool-using large language model (LLM) agents~\cite{gridagent2025}, and a wave of runtime-enforcement systems (SARC~\cite{sarc2026}, AgentSpec~\cite{agentspec2026}, GuardAgent~\cite{guardagent2025}, ShieldAgent~\cite{shieldagent2025}, VeriGuard~\cite{veriguard2025}) makes them more governable. These systems share a rarely-stated assumption: the system starts \emph{inside} the safe set and the gate's job is to keep it there. But infrastructure under autonomous control will sometimes start \emph{outside} it---already in violation, needing a multi-step sequence to recover---and there the enforcement question inverts: the gate must decide whether an action that still leaves the system unsafe nonetheless makes admissible progress toward recovery. This \emph{post-violation recovery admission} problem is structurally different from prevention and largely unaddressed; runtime mediation intercepts unsafe actions yet rarely converts blocked ones into safe completion~\cite{verifiertax2026}. Enforcement is not recovery.

We frame the problem around \emph{admission semantics}: the formal criterion by which a runtime gate decides to apply or reject an LLM-proposed action. Our thesis is that this choice---not the gate's strictness---has first-order effects on whether multi-step recovery is possible. A terminal gate (admit only fully-recovered post-states) is safe but admits no intermediate step, deadlocking every multi-action episode on our mined Gym-ANM~\cite{henry2021gymanm} benchmark. The natural relaxation admits any proposal whose aggregate violation penalty is non-increasing within a small slack---a \emph{scalar admission gate}---and our central finding is that it is structurally insufficient, not for want of threshold tuning.

The reason is the \textbf{scalar projection trap}. In a ReAct loop a rejected proposal is followed by another at the same state, so \emph{the gate is not a filter but a search operator over the agent's proposal stream}; admitting on an aggregate score collapses the multi-component violation geometry and can commit the trajectory to a residual plateau from which recovery may be unreachable, whereas a gate that preserves the per-branch geometry declines the geometry-poor step and keeps the search alive. The gap is representational, which we make precise as a lower bound (Sec.~\ref{sec:method:trap}): structured admission recovers $\mathbf{21/21}$ multi-action episodes where the best scalar slack reaches $9/21$ and terminal $0/21$, with disjoint Wilson 95\% CIs.

\paragraph{Contributions.}
\begin{enumerate}
\item \textbf{The scalar projection trap.} We identify and characterize a failure mode of natural scalar recovery gates: viewed as a search operator over the ReAct proposal stream, a scalar gate can commit the trajectory to a residual plateau because aggregate penalty collapses the per-branch violation geometry the structured order retains (Sec.~\ref{sec:method:trap}). This reframes admission from a safety check into a trajectory-shaping mechanism.
\item \textbf{Structured recovery admission.} We propose \silr{}, a simulator-backed admission layer that shadow-executes each LLM tool proposal and admits it only when a product order over the branch-level violation state is preserved, emitting graded \texttt{PASS}/\texttt{SAFE\_PROGRESS}/\texttt{FAIL} verdicts (Sec.~\ref{sec:method}).
\item \textbf{Empirical validation with a falsification baseline, across families and a second domain.} Terminal admission deadlocks and a scalar threshold sweep cannot close the gap to structured admission; a plateau spectrum and an MPC feasibility check confirm the mechanism is gate-induced (Sec.~\ref{sec:eval:rq1}--\ref{sec:eval:rq3}). The terminal-versus-structured dichotomy holds across three independent model families and reproduces in a second domain (CityLearn): structured admission lets compact open-weight models perform the multi-step recovery that binary and scalar gates cannot (Sec.~\ref{sec:eval:rq5}).
\item \textbf{Trust-boundary containment.} Because admission grounds in deterministic simulation rather than the LLM's input channels, containment is \emph{architectural}, not statistical: a security mini-suite under compromised prompts, observations, and stall instructions yields zero attack-induced worsening or false recovery (Sec.~\ref{sec:eval:rq4}).
\end{enumerate}

\paragraph{Regime and deployment model.}
Our target regime is multi-component concurrent violation---where a scalar penalty aggregates a multi-component severity vector and the projection trap arises---instantiated by Gym-ANM's branch-loading family. \silr{} is an advisory-layer gate at the medium-horizon dispatch loop, where LLM latency is compatible with the pre-dispatch power-flow study. We evaluate this regime in depth and test generalization across three model families and a second domain (CityLearn).
